\textbf{3.\ (Pawn duel).}
Consider the following game on a \(3\times 8\) table.
Each of the two players has three pawns along the short edges.
During the first turn each player may move any of his pawns towards the opposite
edge but not further than by \(3\) cells.
During the later rounds each player may move any of his pawns towards the
opposite edge by any number of cells.
The players cannot move the pawns vertically or diagonally, nor can they jump
across or take pawns of the other player.
The player that has no legal turn loses.
Which player has a winning strategy?
How should he play?

\textbf{Solution.}

For each of the three rows, let
\[
d_i=\text{the number of empty cells between the two pawns in row } i
\]

Initially, in each row the two pawns stand at opposite ends of an \(8\)-cell row, so
\[
d_1=d_2=d_3=6
\]
Thus the initial position is
\[
(6,6,6)
\]

A move in one row simply decreases exactly one of the numbers \(d_i\) by a positive amount. Hence, after the first restricted moves, the game is equivalent to Nim with three heaps of sizes
\[
d_1,d_2,d_3
\]

In Nim, a position is losing if and only if
\[
d_1\oplus d_2\oplus d_3=0,
\]
where \(\oplus\) denotes bitwise XOR.

The initial Nim-sum is
\[
6\oplus 6\oplus 6 = 6\neq 0
\]
Indeed, in binary this is
\[
110_2\oplus 110_2\oplus 110_2 = 110_2
\]

The first player should move one pawn forward by exactly \(2\) cells. This changes one gap from \(6\) to \(4\), so the position becomes
\[
(4,6,6).
\]

Now
\[
4\oplus 6\oplus 6 = 4
\]

If the second player wanted to make the Nim sum zero immediately, he would have to change one heap \(h\) into
\[
h' = h\oplus 4
\]

The possible changes for \(h\) are :

\(h=4\):
\[
4\to 4\oplus 4 = 0,
\]
or for \(h=6\):
\[
6\to 6\oplus 4 = 2
\]

Both moves require decreasing a gap by \(4\). But on the second player's first move, he is allowed to move a pawn by at most \(3\) cells. Therefore, the second player cannot move to a zero Nim-sum position.

After that, the first player is no longer restricted. From then on he uses the standard Nim strategy: after every move of the opponent, make
\[
d_1\oplus d_2\oplus d_3=0
\]

For example, after the first move
\[
(6,6,6)\to(4,6,6),
\]
the second player's first move can produce, up to permutation of rows, only one of the following positions:
\[
(3,6,6),\quad (2,6,6),\quad (1,6,6),
\]
or
\[
(4,5,6),\quad (4,4,6),\quad (4,3,6)
\]

The first player answers as follows:
\[
(3,6,6)\to(0,6,6),
\]
\[
(2,6,6)\to(0,6,6),
\]
\[
(1,6,6)\to(0,6,6),
\]
\[
(4,5,6)\to(3,5,6),
\]
\[
(4,4,6)\to(4,4,0),
\]
\[
(4,3,6)\to(4,2,6)
\]

In each case the resulting Nim-sum is zero.

Thus the first player can always give the second player a position satisfying
\[
d_1\oplus d_2\oplus d_3=0.
\]

The game eventually ends.
Since the first player always leaves the second player a zero Nim sum position, the second player is the one who eventually has no legal move.

\smallbreak
Therefore, the first player has a winning strategy.

\[
\text{The first player wins.}
\]

His strategy is:

\[
\text{First move: move any pawn forward by exactly }2\text{ cells.}
\]

After that, always move so that
\[
d_1\oplus d_2\oplus d_3=0
\]
